% !TEX root = ../main.tex
\chapter*{LỜI MỞ ĐẦU}
\addcontentsline{toc}{chapter}{LỜI MỞ ĐẦU}

Trong bối cảnh bùng nổ của cuộc Cách mạng Công nghiệp lần thứ tư, khối lượng các công bố khoa học quốc tế tăng trưởng theo cấp số nhân mỗi năm. Các học giả, giảng viên và nhà nghiên cứu phải đối mặt với áp lực khổng lồ trong việc sàng lọc, tiếp thu và tổng hợp tri thức từ hàng nghìn bài báo chuyên ngành. Mặc dù các mô hình ngôn ngữ lớn (Large Language Models -- LLM) đã đạt được những bước tiến vượt bậc trong xử lý ngôn ngữ tự nhiên, việc ứng dụng các chatbot hội thoại phổ thông vào nghiên cứu học thuật vẫn vấp phải rào cản nghiêm trọng mang tên hiện tượng ảo giác trích dẫn (Citation Hallucination). Các mô hình thường tự tạo ra các trích dẫn không có thật, hoặc trả lời một cách khái quát mà không thể đối chiếu chính xác tới từng số liệu, bảng biểu hay trang tài liệu gốc.

Hơn thế nữa, quy trình nghiên cứu khoa học của một học giả vốn là một chuỗi tác vụ liên tục và khép kín: từ việc khám phá bài báo mới, đọc hiểu chuyên sâu với sự hỗ trợ của AI, trích xuất mạng lưới phương pháp -- bộ dữ liệu -- kết quả thực nghiệm, cho đến bước soạn thảo bản thảo khoa học bằng \LaTeX. Hiện nay, các giải pháp hỗ trợ vẫn còn mang tính chất rời rạc, thiếu đồng bộ: công cụ đọc PDF không giữ được bố cục 2 cột và không thể định vị hộp tọa độ; thư viện quản lý tài liệu thiếu khả năng suy luận liên bài báo; trong khi các trình soạn thảo \LaTeX\ trực tuyến đòi hỏi người viết phải tự sửa hàng loạt lỗi cú pháp biên dịch phức tạp.

Xuất phát từ yêu cầu cấp thiết của thực tiễn học thuật và mong muốn ứng dụng các tiến bộ mới nhất của trí tuệ nhân tạo vào hỗ trợ nghiên cứu, đề tài \textit{''Xây dựng nền tảng AI hỗ trợ nghiên cứu và viết bài báo khoa học VinUni dựa trên Agentic RAG, truy xuất tri thức học thuật và kiểm chứng trích dẫn''} được lựa chọn thực hiện. Đồ án hướng tới việc thiết kế và phát triển một hệ sinh thái nghiên cứu khoa học chuyên sâu, khép kín 4 không gian làm việc, cung cấp bằng chứng tọa độ chính xác cho từng câu trả lời của AI và hỗ trợ tự động hóa chu trình biên soạn bản thảo học thuật đạt chuẩn quốc tế.
